% RESEARCH_DRAFT_SHA256: f6dbc0ce88fea8ba863e52efa4b1eeaac7fe1295cdea2d5f61763e9cc533b58b
% Layout and source-reference wrappers expanded for the web renderer.
\documentclass[11pt]{article}
\usepackage[margin=25mm]{geometry}
\usepackage{fontspec,unicode-math}
\setmainfont{Latin Modern Roman}
\setmathfont{Latin Modern Math}
\usepackage{amsmath,amssymb,mathtools,xurl,hyperref}
\hypersetup{colorlinks=true,urlcolor=blue,linkcolor=blue}
\setcounter{secnumdepth}{0}
\setlength{\parindent}{0pt}
\setlength{\parskip}{5pt}
\setlength{\emergencystretch}{3em}
\newcommand{\sref}[2]{\hyperlink{source:#1:#2}{[S#2]}}
\newcommand{\cataloguescope}[1]{\paragraph{Source record.}#1}
\newcommand{\floor}[1]{\left\lfloor#1\right\rfloor}
\begin{document}
\section*{Research-attempt provenance}
This standalone research attempt retains the complete original definition body below,
including its historical source and preparation notes. Those notes describe the
earlier catalogue preparation; the new attempt and its verification are identified
separately after that body. No original ID, catalogue status, or source statement is changed.

\section{A versioned capability contract for PreTeXt programming languages}\label{statement}

\subsection{Definitions and mathematical statement}

\paragraph{Shared notation.}Standard number systems are $\mathbb N,\mathbb Z,\mathbb Q,\mathbb R,\mathbb C$. Specialized notation and index ranges follow the cited source and the definitions below.


A programming language specifies rules for expressing computations. An embedded execution environment supports particular languages and libraries with defined inputs and outputs. Capabilities can include runtime access, interactivity, available packages, and resource constraints.
\paragraph{Record scope and qualifications.}
The catalogue title is retained as metadata. Named certificates, contracts, bounds or frameworks in that title are not asserted as results of the original source extract.
The platform, execution environment, and relevant capability criteria are missing; no current availability list is asserted.

Catalogue screening limitation: PreTeXt software design, documentation, or publication feature request; no mathematical problem is specified by the supplied passage.

The source snapshot labels this record partially solved; this is catalogue metadata rather than independent verification of a mathematical result or the current completion of a project.
\paragraph{Statement recorded in the supplied source.}
Programming Languages: Which are available, what are their capabilities? \href{https://github.com/PreTeXtBook/community/wiki/Moderated-Problem-Session}{[S2]}
\paragraph{Source status.}
The supplied catalogue labels this source record partially solved; that metadata was checked against the supplied snapshot on 2026-10-11. The label is preserved for provenance and is not an independent verification of present open status or mathematical correctness.
\subsection{Short English statement}
Catalogue available programming languages and their capabilities in the intended document environment.
\subsection{Sources}
\begin{itemize}
\item[S1] Ulam AI and the UnsolvedMath Contributors, UnsolvedMath: A Curated Collection of Open Mathematics Problems, catalogue entry. Catalogue curation and metadata: CC BY 4.0. The source extract is retained; definitions, English summaries and qualifications are editorial additions.
\url{https://huggingface.co/datasets/ulamai/UnsolvedMath}.
\item[S2] AIM Workshop Problem Lists, source linked by the catalogue.
\url{https://github.com/PreTeXtBook/community/wiki/Moderated-Problem-Session}.
\end{itemize}
\paragraph{Preparation and review.}
The original source and UnsolvedMath contributors retain attribution for the source extract. Codex (GPT-6; exact serving revision not exposed) assisted with definitions, English summaries, scope notes and formatting. No new human mathematical authorship or independently verified solution is asserted. Catalogue identity checked against the complete supplied registry and current repository. Distinct upstream records, including repeated questions and broad or damaged extracts, are retained for complete coverage of the source CS category. A repeated record is not a new independent conjecture. Complete mathematical-equivalence and current-literature review has not been performed. Preparation and final review follow contribution prompts C and D, with source gaps disclosed. Structural and rendering checks do not establish a new theorem or open status.
\label{end:statement}

\section{Literature and context check}
The primary community-session source was consulted. The record is a broad software-design request, and no mathematical theorem follows from its editorial catalogue title. The finite model below is explicitly auxiliary; it does not certify an actual implementation or current feature status. The consulted PreTeXt schema snapshot admits a textual language attribute on program elements, among other execution-related attributes. Syntactic admission of that string does not establish an available runtime or its permissions.
\section{Research attempt: A language name is not an execution capability certificate}
Consider two abstract runtimes that accept the same language name. Runtime $A$ evaluates arithmetic expressions only; runtime $B$ additionally reads a supplied data file. The arithmetic task $2+2$ succeeds in both, but a task requiring the file's first entry is undefined in $A$. A capability table with one Boolean cell saying that the language is available identifies these runtimes even though a required task distinguishes them. Thus language identity and task capability are different data.

For a finite declared set of tasks $T$, represent evidence by tuples $(\text{runtime},\text{version},\text{environment},t,\text{outcome})$. A task can be certified for a selected tuple only if that task's required operations and expected observations were tested under those same coordinates. Combining arithmetic success from one environment with file-access success from another need not certify their combined task: neither individual environment may support both. The counterexample uses environment 1 with arithmetic only and environment 2 with file access only, while the combined task needs both.

In a finite operation-set model, let a task require a set $Q$ of operations and let a runtime supply $C$. The necessary inclusion is $Q\subseteq C$. If the model defines success exactly as availability of all required operations with compatible semantics, the inclusion is also sufficient by that definition. Actual runtime success has additional constraints such as input formats and resource bounds; those cannot be proved from a set of operation names. The useful conclusion is a precise failure of an unqualified language-support matrix, not a list of current supported programming languages.
\section{Verification and adversarial audit}
The construction and counterexample are checked in the written finite model. No actual PreTeXt build, browser, EPUB reader, or publishing service was tested. All runtime and identity assumptions used by the proof are stated explicitly.
\paragraph{Audit corrections.} A schema attribute is not promoted to execution evidence. Task conjunction is evaluated within one runtime/version/environment rather than pooled across incompatible observations.
\paragraph{External computation.} None; symbolic arguments and boundary checks only.
\section{Scope and remaining gaps}
This is a restricted route investigation. The full feature requires implementation, supported-version decisions, conformance tests, and user-facing evaluation that were not performed here. The original agenda remains unresolved, and no novelty or completed software feature is claimed.
\paragraph{Integration review.} Codex (GPT-6; exact serving revision not exposed) reviewed the written arguments and expanded layout and source-reference wrappers. Repository metadata and the production web renderer were checked using Python 3.13 and JavaScript with MathJax 3.2.2. This model-assisted review is not independent human proof verification; no PDF compilation or formal proof checking is asserted.

\section{Completion Estimate}
\noindent\textbf{COMPLETION: 0\%}
The calculations below the preserved source statement investigate a concrete route,
but establish no substantiated new progress toward the entire original target.
Known special cases and failed sufficient conditions are not credited as a solution.
\section{References}
The original source references and catalogue attribution are preserved in the
Sources section above. The following URLs were actually consulted for this attempt;
retention of an additional original reference does not claim fresh consultation of it.
\begin{itemize}
\item \url{https://github.com/PreTeXtBook/community/wiki/Moderated-Problem-Session} (consulted 2026-10-11). 
\item \url{https://github.com/PreTeXtBook/pretext/blob/master/schema/pretext.rnc} (consulted 2026-10-11). 
\end{itemize}
\section{Model and contributor disclosure}
OpenAI gpt-6-astra, requested reasoning configuration ultra, generated this research
attempt on 2026-10-11. The exact serving revision is not exposed. Codex GPT-6
performed orchestration and repository-format review; its exact serving revision
is also not exposed. The mathematical route, boundary checks and final adversarial
self-audit followed contribution prompts A then B in
\texttt{docs/contribution\_prompts.md}, after reading \texttt{CONTRIBUTING.md}.
Python 3.14.2 assembled the standalone documents; any mathematical computation is
specified in Verification. No theorem prover was used. No human mathematical
authorship, substantial human proof correction, or independent human proof checking
is asserted. Model-assisted review, source retrieval, and structural checks do not
establish correctness or novelty. Full private reasoning transcripts are unavailable
as reproducibility artifacts; the complete arguments, inputs and outcomes actually
used are stated here or in the companion batch audit.
\section{Final status}
Unresolved as a new research contribution. This is a concrete route investigation,
not a claim of new complete resolution or a current-status certification. A known
answer reproduced for context does not become a new contribution. The catalogue's
source-status label is preserved; this attempt status does not assert that a
source-labelled solved question is open.
\end{document}
